
% Authoritative LaTeX manuscript; Markdown is a generated compatibility projection.

Provenance models record what happened. Documentation models represent the human-facing claims that should remain synchronized with those facts. In review, they expose whether required records exist and whether derived documentation still reflects its source.

\textbf{Engineering question.} What happened, and which human-facing account of it must remain true?

\textbf{Representation.}

% mage-code-language: text
\begin{verbatim}
Action
    id
    artifact
    operation
    actor
    mechanism
    timestamp
    result
DerivedDocumentation
    source_fact
    rendered_claim
\end{verbatim}

Where documentation can be derived from structured facts, generate it from those facts rather than maintaining an independent copy.

\autoref{fig-g6-facts-before-prose} puts the fact before the prose.

\begin{figure}[htbp]
\centering
\includegraphics[width=\linewidth]{book/latex/assets/appendix-g-6-facts-before-prose-525483293a.pdf}
\caption{\emph{Facts before prose.} A structured provenance record supports mechanical validation and re-derivable documentation; independently maintained prose lacks that correspondence.}
\label{fig-g6-facts-before-prose}
\end{figure}

\textbf{Property.} Representative properties:

\begin{itemize}
\tightlist
\item
  \textbf{Every relevant mutation leaves a provenance record.}
\item
  \textbf{Generated documentation reflects the current structured facts} from which it is derived.
\item
  \textbf{References in generated documentation resolve to existing identities.}
\end{itemize}

\textbf{Enforcement and correspondence.} Structured provenance can be validated mechanically, and generated documentation can be re-derived from its source facts. Free prose generally cannot be checked for semantic equivalence by the same machinery and therefore remains subject to human review.
